\section{Introduction}
\label{sec:introduction}

Accumulation is a central operation in mathematics and computing --- from
prefix sums in algorithms to integral transforms in signal processing. In
functional programming, accumulation often appears as a fold or scan, but
such constructs are rarely defined from first principles in a formally
verified setting.

In this article, we define a recursive discrete integral over finite
integer lists and verify its cumulative-sum, difference-recovery,
monotonicity, and representation-agreement properties using Scala
Stainless. The recursive lookup and accumulated-list representations agree
pointwise and in length, and consecutive differences recover the
corresponding input values.

This article verifies:

\begin{itemize}
  \item Core integral properties: head value, cumulative sum, incremental
    change, final sum, strictly increasing, gaps positivity
    --- Sections~4.1--4.6.
  \item Implementation consistency: element/acc/delta/last/size agreement
    between the recursive and accumulated representations
    --- Sections~5.2--5.5.
\end{itemize}

\subsection*{Related work}

The cumulative-sum construction is the list instance of a prefix scan or
accumulation. In Rocq/Coq, the standard list library defines
\texttt{fold\_left} and proves its composition across concatenation; its
natural-number list library also defines list sum as a fold and proves sum
over concatenation \cite{rocqliststdlib}. Those formally checked results
give a useful established setting for recursive accumulation.

The present article develops that setting for a recursive \texttt{BigInt}
integral in Scala Stainless. Its focus is the agreement of two concrete
representations --- recursive lookup and an accumulated list --- and the
accompanying cumulative-sum, difference-recovery, length, and monotonicity
properties. The citation places these proofs in a broader formal treatment
of list accumulation; it does not replace the specific
representation-agreement results verified here.
